\section{From large marginals to a finite degree sum}
\label{sec:lift}

Proposition~\ref{prop:marginals} gives one partition for which the
complementary image laws have small total error. We first express those
image laws as coset laws and condition on one event that bounds all their
coset densities. We then turn these bounds into a finite sum of squared
degrees of representations of the block groups. The next section will
estimate that sum by counting tableaux.

Throughout this section, \(U_G\) denotes the uniform probability measure on
a finite group \(G\). For a subgroup \(H\), we use the standard convention
that \(gH\) is a \emph{left coset}. Its elements are obtained from \(g\) by
multiplication on the right by elements of \(H\). If \(\mu\) is a probability
measure on \(G\), write \(\mu_H\) for its pushforward to the set \(G/H\) of
left cosets. Thus
\[
  \mu_H(gH)=\mu(gH),
  \qquad (U_G)_H(gH)=\frac{|H|}{|G|}.
\]

For the shuffle application, take $X=V$, set $T=1025d$ and
$\mu=q_d^{*T}$, and use the deterministic partition
$M_1,\ldots,M_k$ supplied by Proposition~\ref{prop:marginals},
with $k=64$ and $|M_a|=n/k$. Let $H_a$ permute $M_a$ and fix
$X\setminus M_a$ pointwise.
A permutation maps each reference label to its current position. Two maps
\(g,g'\) agree on \(X\setminus M_a\) exactly when
\(g^{-1}g'\in H_a\), or equivalently \(g'\in gH_a\). Thus the ordered
images of the cards outside \(M_a\) identify the coset \(gH_a\)
bijectively. Uniform measure on permutations induces uniform measure on
these ordered distinct images. The marginal distances in
Proposition~\ref{prop:marginals} are consequently exactly the coset
distances \(\delta_a\) for the trimming lemma below, and
\(\sum_a\delta_a=o(1)\).

\begin{lemma}[Simultaneous trimming]
\label{lem:trim}
Let \(H_1,\ldots,H_k\) be subgroups of a finite group \(G\), and let
\(\mu\) be a probability measure on \(G\). Put
\[
  \delta_a=\|\mu_{H_a}-(U_G)_{H_a}\|_{\TV},
  \qquad \delta=\sum_{a=1}^k\delta_a.
\]
If \(\delta<1/2\), there is a set \(E\subseteq G\), of removed mass
\(\varepsilon=\mu(G\setminus E)\le 2\delta\), such that the conditional
probability measure \(\bar\mu=\mu(\,\cdot\mid E)\) satisfies
\begin{equation}
\label{eq:lift:trim}
  \|\bar\mu-\mu\|_{\TV}=\varepsilon,
  \qquad
  \bar\mu(gH_a)
  \le \frac{2}{1-\varepsilon}\frac{|H_a|}{|G|}
  \quad(g\in G,\ 1\le a\le k).
\end{equation}
In particular, if \(\delta\le 1/4\), the last density bound is at most
\(4|H_a|/|G|\).
\end{lemma}

\begin{proof}
For each \(a\), call a coset \(C\in G/H_a\) bad if
\(\mu(C)>2U_G(C)\), and let \(B_a\) be the union of its bad cosets.
The identity expressing total variation as the sum of positive excesses
gives
\[
  \mu(B_a)
  \le 2\sum_{\substack{C\in G/H_a\\ \mu(C)>2U_G(C)}}
                 \bigl(\mu(C)-U_G(C)\bigr)
  \le 2\delta_a.
\]
Take \(E=G\setminus\bigcup_a B_a\). The union bound yields
\(\varepsilon\le 2\delta<1\), so conditioning is well defined. On \(E\),
the mass of \(\bar\mu\) exceeds that of \(\mu\) by total amount
\(\varepsilon\); on the complement, its mass is zero. This proves the
total-variation equality in \eqref{eq:lift:trim}.

Fix a coset \(C\in G/H_a\). If it is bad for this subgroup, then
\(C\cap E=\varnothing\). Otherwise,
\[
  \bar\mu(C)
  =\frac{\mu(C\cap E)}{1-\varepsilon}
  \le \frac{\mu(C)}{1-\varepsilon}
  \le \frac{2U_G(C)}{1-\varepsilon}.
\]
The same estimate therefore holds for every coset in every one of the
\(k\) systems. If \(\delta\le 1/4\), then
\(1-\varepsilon\ge 1/2\), giving the last assertion.
\end{proof}

We now convert the coset bounds into control of averages of representation
matrices. For a mass function \(a:G\to\mathbb C\) and a unitary
representation \(\rho\), define its Fourier matrix by
\[
  \widehat a(\rho)=\sum_{g\in G}a(g)\rho(g).
\]
Probability and subprobability measures are identified with their mass
functions. For a finite group \(H\), write \(\widehat H\) for the set of
equivalence classes of its complex irreducible representations.

When an irreducible representation of \(G\) is restricted to the product
of the block groups, each irreducible summand is a tensor product of one
representation from each group. If the full representation has dimension
\(D\), each such tensor product has dimension at most \(D\), so at least
one factor has degree at most \(D^{1/k}\). The following proposition
measures the cost of these small factors by a finite sum. Its proof shows
why each equivalence class of a subgroup representation contributes only
the square of its degree even when it occurs with large multiplicity.

\begin{proposition}[Coset densities and squared subgroup degrees]
\label{prop:lift:orbit-sum}
Let a finite set \(X\) be partitioned into \(k\ge1\) nonempty sets
\(M_1,\ldots,M_k\), and put \(G=\Sym(X)\). Let \(H_a\) be the subgroup
permuting \(M_a\) and fixing \(X\setminus M_a\) pointwise. Suppose that a
probability measure \(\nu\) on \(G\) satisfies, for some \(B\ge 1\),
\begin{equation}
\label{eq:lift:coset-cap}
  \nu(gH_a)\le B\frac{|H_a|}{|G|}
  \qquad(g\in G,\ 1\le a\le k).
\end{equation}
Then every irreducible complex unitary representation \(\rho\) of \(G\), of
dimension \(D\), satisfies
\begin{equation}
\label{eq:lift:orbit-sum}
  \|\widehat\nu(\rho)\|_{\op}^2
  \le \frac{B}{D}\sum_{a=1}^k
       \sum_{\substack{\pi\in\widehat H_a\\\dim\pi\le D^{1/k}}}
                    (\dim\pi)^2.
\end{equation}
The block sizes need not be equal.
\end{proposition}

\begin{proof}
Because the blocks are disjoint, their subgroups form the direct product
\(H=H_1\times\cdots\times H_k\) inside \(G\). By complete reducibility
and the irreducible-representation theorem for finite direct products,
the restriction of \(\rho\) to \(H\) is an orthogonal sum of isotypic
spaces of the form
\begin{equation}
\label{eq:lift:restriction}
  (V_{\pi_1}\otimes\cdots\otimes V_{\pi_k})
       \otimes \mathcal A_{\boldsymbol\pi}.
\end{equation}
Here \(\pi_a\) is an irreducible representation of \(H_a\), and
\(\mathcal A_{\boldsymbol\pi}\) is a multiplicity space on which \(H\)
acts trivially. These are the usual characteristic-zero representation
facts; see, for example, \cite{sagan,etingof}.

For every type \(\boldsymbol\pi\) that occurs, its tensor-product
dimension is at most \(D\). Consequently some index \(a\) satisfies
\(\dim\pi_a\le D^{1/k}\). Assign each entire isotypic space in
\eqref{eq:lift:restriction} to one such index, choosing the least one
in case of a tie. Grouping the assigned spaces yields an orthogonal
decomposition
\[
  V_\rho=E_1\oplus\cdots\oplus E_k.
\]
Each \(E_a\) is invariant under \(H\), and every irreducible of \(H_a\)
appearing in \(E_a\) has dimension at most \(D^{1/k}\). The multiplicities
in these restrictions are unrestricted.

For this fixed representation, put
\[
  r_a=\sum_{\substack{\pi\in\widehat H_a\\\dim\pi\le D^{1/k}}}
                (\dim\pi)^2.
\]
Fix \(u_a\in E_a\), and set
\[
  W_a=\Span\{\rho(h)u_a:h\in H_a\}.
\]
We first bound this single-vector orbit span. Regroup all copies of each
\(H_a\)-irreducible, including copies coming from different \(H\)-types,
into a single \(H_a\)-isotypic space. In such an isotypic
space \(V_\pi\otimes\mathcal A_\pi\), with \(b=\dim\pi\), the group
acts by \(\pi(h)\otimes\id\). If \(v\) is the projection of \(u_a\)
to this space, its orbit span is contained in the image of the linear map
\[
  \operatorname{End}(V_\pi)\longrightarrow
          V_\pi\otimes\mathcal A_\pi,
  \qquad Q\longmapsto (Q\otimes\id)v.
\]
That image has dimension at most \(b^2\), independently of
\(\dim\mathcal A_\pi\). Each inequivalent \(H_a\)-type contributes only
once. Summing over the different isotypic spaces gives the finite bound
\begin{equation}
  \dim W_a\le r_a
  =\sum_{\substack{\pi\in\widehat H_a\\\dim\pi\le D^{1/k}}}
                         (\dim\pi)^2.
\label{eq:lift:orbit-dimension}
\end{equation}

Let \(\Pi_a\) be the orthogonal projection onto \(W_a\). The subspace
\(W_a\) is \(H_a\)-invariant, and therefore
\(\rho(gh)W_a=\rho(g)W_a\) for \(h\in H_a\). For each vector \(z\),
the nonnegative function
\[
  g\longmapsto
       \|\Pi_{\rho(g)W_a}z\|^2
\]
is thus constant on each left coset \(gH_a\). The coset bound
\eqref{eq:lift:coset-cap} implies
\begin{equation}
\label{eq:lift:projection-comparison}
  \E_{g\sim\nu}\|\Pi_{\rho(g)W_a}z\|^2
  \le B\E_{g\sim U_G}\|\Pi_{\rho(g)W_a}z\|^2.
\end{equation}
The uniform average of these projections is
\[
  Q_a=\frac1{|G|}\sum_{g\in G}\rho(g)\Pi_a\rho(g)^*.
\]
It commutes with every \(\rho(g)\). Schur's lemma makes it a scalar
multiple of \(\id\), and its trace is \(\dim W_a\). Hence
\begin{equation}
\label{eq:lift:schur-average}
  Q_a=\frac{\dim W_a}{D}\id,
  \qquad
  \E_{g\sim U_G}\|\Pi_{\rho(g)W_a}z\|^2
       =\frac{\dim W_a}{D}\|z\|^2.
\end{equation}

Since \(\rho(g)u_a\in\rho(g)W_a\), Cauchy--Schwarz, first for
expectation and then in the representation space, gives
\begin{align}
  |\langle z,\widehat\nu(\rho)u_a\rangle|^2
  &\le \E_{g\sim\nu}|\langle z,\rho(g)u_a\rangle|^2 \notag\\
  &\le \|u_a\|^2
           \E_{g\sim\nu}\|\Pi_{\rho(g)W_a}z\|^2 \notag\\
  &\le \frac{Br_a}{D}\|u_a\|^2\|z\|^2,
\label{eq:lift:component-bound}
\end{align}
where the last line uses
\eqref{eq:lift:orbit-dimension}--\eqref{eq:lift:schur-average}.
For an arbitrary \(u=\sum_a u_a\), sum the square-root form of this
estimate and apply Cauchy--Schwarz in the index \(a\):
\begin{align*}
 |\langle z,\widehat\nu(\rho)u\rangle|
 &\le \sqrt{\frac{B}{D}}\,\|z\|
           \sum_{a=1}^k\sqrt{r_a}\,\|u_a\|\\
 &\le \sqrt{\frac{B}{D}\sum_{a=1}^k r_a}\,\|z\|
           \left(\sum_{a=1}^k\|u_a\|^2\right)^{1/2}\\
 &=\sqrt{\frac{B}{D}\sum_{a=1}^k r_a}\,\|z\|\,\|u\|.
\end{align*}
The last equality uses the orthogonality of the \(E_a\). Taking the
supremum over unit vectors \(u,z\) and squaring proves
\eqref{eq:lift:orbit-sum}. The argument includes zero components and
one-dimensional representations; in dimension one the displayed sum
does not provide decay.
\end{proof}
